\section{Certified short-chain thermal moments}\label{app:thermal}

This appendix proves the two short-chain estimates used to initialize the
boundary argument. Throughout, $h=3/5$, $a=700741/500000$, and $b_0=49/4$.
All terminating decimals below denote exact rational numbers.

\begin{proposition}\label{appth:bounds}
Let $\lambda_0$ be the positive leading eigenvalue of $X_{b_0}$, and let
$c_0$ be the squared norm of the projection of $v_{b_0}$ onto its eigenspace.
For
\[
 P_*(x)=x^3-.14x^2-.18x,
\]
one has
\begin{equation}\label{appth:conclusions}
 Y_6(b_0)<.027,
 \qquad c_0P_*(\lambda_0)^2<.003745.
\end{equation}
\end{proposition}

The first bound controls the residual contribution to $Y_m(b_0)$ in
Lemma~\ref{in:seedbounds}; the second controls the leading contribution.
We first reduce both bounds to fifteen short-chain partition functions,
then enclose these traces using matrices of order at most $3^6=729$.

\subsection{Reduction to fifteen thermal traces}

For $2\le n\le6$ and $g\in\{1,-1,0\}$, define
\begin{align}
 H_{n,g}
 &=\sum_{j=0}^{n-2}\mathbf S_j\mathbin\cdot\mathbf S_{j+1}
       -\frac35S_0^z+\frac35R_g,
 &
 R_g&=
 \begin{cases}
    gS_{n-1}^z,&g=1,-1,\\
    S_{n-1}^x,&g=0,
 \end{cases}
 \label{appth:H}\\
 Z_{n,g}&=\Tr\exp[-b_0(H_{n,g}+an\Id)].\nonumber
\end{align}
Thus $Z_{n,1}=Y_n(b_0)$ and $Z_{n,-1}=Y_n^{\rm s}(b_0)$.
The first bound in Proposition~\ref{appth:bounds} concerns $Z_{6,1}$.
For the second, average the right-field direction over the rotation group.

Let $\mathcal G$ be the rotation group generated by the diagonal half-turns and the
cyclic permutation of the coordinate axes, and let
\[
 \Pi=\frac1{|\mathcal G|}\sum_{\gamma\in\mathcal G}U_\gamma,
 \qquad I(n)=\langle v_{b_0},X_{b_0}^n\Pi v_{b_0}\rangle.
\]
The group has order $12$, and the orbit of $z$ is the six signed coordinate
axes, each with the same multiplicity. The boundary transfer identity
therefore writes $I(n)$ as the average of the six corresponding right-field
partition functions. Global rotations about $z$ identify the four
perpendicular orientations. Hence
\begin{equation}\label{appth:average}
 I(n)=\frac{Z_{n,1}+Z_{n,-1}+4Z_{n,0}}6.
\end{equation}
The averaging operator $\Pi$ is an orthogonal projection commuting with
$X_{b_0}$. By Proposition~\ref{per:concentration}, $\mathcal G$ fixes the
leading line pointwise. Therefore
\begin{align}
 c_0P_*(\lambda_0)^2
 &\le\|P_*(X_{b_0})\Pi v_{b_0}\|^2\nonumber\\
 &=\langle v_{b_0},P_*(X_{b_0})^2\Pi v_{b_0}\rangle\nonumber\\
 &=I(6)-.28I(5)-.3404I(4)+.0504I(3)+.0324I(2).
 \label{appth:polynomial}
\end{align}
In particular, no positivity of the individual odd powers of $X_{b_0}$
is being assumed.

It therefore suffices to enclose $Z_{n,g}$ for $2\le n\le6$ and
$g\in\{1,-1,0\}$. The remaining subsections provide these enclosures
and evaluate \eqref{appth:polynomial}.

\subsection{An integer matrix for the thermal approximation}

Set
\begin{equation}\label{appth:C}
 C_n=3\left\lfloor\frac{n-1}{2}\right\rfloor
          +2\bigl((n-1)\bmod2\bigr)+\frac65,
 \qquad A=\Id-\frac{H_{n,g}+C_n\Id}{8}.
\end{equation}

The operator $A$ is a self-adjoint contraction. Here is a direct verification
with the spin normalization in the theorem. A single bond has eigenvalues
$-2,-1,1$. For two consecutive bonds, couple the two outside spins to spin
$t\in\{0,1,2\}$ and then couple this with the middle spin to total spin $j$.
The operator
\[
 \mathbf S_1\mathbin\cdot\mathbf S_2+
 \mathbf S_2\mathbin\cdot\mathbf S_3
   =\mathbf S_2\mathbin\cdot(\mathbf S_1+\mathbf S_3)
\]
has eigenvalues $[j(j+1)-t(t+1)-2]/2$, with
$|t-1|\le j\le t+1$. Its minimum is $-3$.
Pairing consecutive bonds, with one single bond left over if necessary,
and using the total boundary-field norm $6/5$ gives
\[
 0\le H_{n,g}+C_n\Id
    \le\left(n-1+\frac65+C_n\right)\Id\le16\Id.
\]
Adding the lower bounds for the paired bonds is valid even where the
three-site supports overlap. In order $n=2,\ldots,6$, the displayed upper
bounds are $5.4,7.4,10.4,12.4,15.4$. Thus $-\Id\le A\le\Id$;
positivity of $A$ itself is not required.

Index the spin-weight basis by words $w=(w_0,\ldots,w_{n-1})$ in
$\{-1,0,1\}^n$, and let
\[
 p_w=\left(\sum_{j=0}^{n-1}w_j\right)\bmod2\in\{0,1\},
 \qquad J_{ww}=\sqrt{1+p_w},\qquad D=80JAJ^{-1}.
\]
The matrix $D$ is integral. Indeed, its diagonal is
\begin{equation}\label{appth:diagonal}
 \eta(w)=80-10C_n-10\sum_{j=0}^{n-2}w_jw_{j+1}
                  +6w_0-6g w_{n-1}.
\end{equation}
Each allowed bond hop $w\leftrightarrow
w+\sigma(\mathbf e_j-\mathbf e_{j+1})$, $\sigma=\pm1$, has entry $-10$:
the corresponding entry of the spin-one bond is $1$, and such a hop
preserves $p_w$. When $g=0$, an allowed end hop from column
$r=w+\sigma\mathbf e_{n-1}$ to row $w$ has entry $-3(1+p_w)$.
To see this coefficient, the end-field entry before similarity
is $3\sqrt2/10$, and the two parities are opposite, so
\[
 -10\frac{3\sqrt2}{10}
       \sqrt{\frac{1+p_w}{1+p_r}}=-3(1+p_w).
\]
These are all the entries of $D$.

\subsection{The exact recurrence}

The recurrence approximates the half-temperature exponential
\[
 E=\exp[-b_0(H_{n,g}+C_n\Id)/2]=e^{-49}\exp(49A).
\]
For a finite matrix $M$, write $\|M\|_{\rm F}^2=\Tr(M^*M)$ for
its squared Frobenius norm. Since $E$ is self-adjoint, $\|E\|_{\rm F}^2$
is the thermal trace $\Tr\exp[-b_0(H_{n,g}+C_n\Id)]$. We approximate
$E$ by a normalized Poisson polynomial in the contraction $A$;
the similarity $D=80JAJ^{-1}$ permits an integer Horner recurrence.

Put
\begin{equation}\label{appth:coefficients}
 K=192,\qquad Q=2^{50},\qquad
 s_i=\frac{49^iK!}{i!},\qquad S_K=\sum_{i=0}^Ks_i,
 \qquad t_i=\left\lfloor\frac{Qs_i}{S_K}\right\rfloor.
\end{equation}
Starting with the zero matrix $\mathcal X_{K+1}$, compute
\begin{equation}\label{appth:horner}
 \mathcal X_i=
       \left\lfloor\frac{D\mathcal X_{i+1}}{80}\right\rfloor+t_i\Id,
 \qquad i=K,K-1,\ldots,0.
\end{equation}
Matrix floors are entrywise, and every floor, including one with a negative
argument, is toward $-\infty$. Define
\begin{equation}\label{appth:W}
 \widehat E=J^{-1}\frac{\mathcal X_0}{Q}J,
 \qquad W_{n,g}=\|\widehat E\|_{\rm F}^2.
\end{equation}

For completeness, the following scalar version specifies the recurrence
without matrix construction. Number words, if desired, by
$\sum_j(w_j+1)3^j$, and interpret a lookup outside $\{-1,0,1\}^n$ as zero.
Initialize $x(w,r)=0$ for all pairs of words. At stage $i=K,\ldots,0$, use
the current array to form
\begin{align}
 z(w,r)={}&\eta(w)x(w,r)
       -10\sum_{j=0}^{n-2}\sum_{\sigma=\pm1}
             x\bigl(w+\sigma(\mathbf e_j-\mathbf e_{j+1}),r\bigr)
       \nonumber\\
 &\hspace{5mm}
       -\mathbf1_{\{g=0\}}\,3(1+p_w)
             \sum_{\sigma=\pm1}x(w+\sigma\mathbf e_{n-1},r),
       \label{appth:scalar}\\
 x_{\rm new}(w,r)={}&t_i\mathbf1_{\{w=r\}}
                    +\left\lfloor\frac{z(w,r)}{80}\right\rfloor.
       \nonumber
\end{align}
Replace all entries simultaneously. At the end, the required table entry is
the integer quotient
\begin{equation}\label{appth:quotient}
 k_{n,g}:=\left\lfloor10^{16}W_{n,g}\right\rfloor
 =\left\lfloor
 \frac{10^{16}\displaystyle\sum_{w,r}
                 x(w,r)^2(2-p_w)(1+p_r)}{2Q^2}
 \right\rfloor.
\end{equation}
The weight follows from
$(1+p_r)/(1+p_w)=(2-p_w)(1+p_r)/2$.
All operations in \eqref{appth:coefficients}, \eqref{appth:scalar}, and
\eqref{appth:quotient} are integer operations; they may be performed with
arbitrary-precision integers.

The resulting fifteen entries are
\begin{equation}\label{appth:integer-table}
\begin{array}{c|rrr}
 n\backslash g&1&-1&0\\\hline
 2&10185296075298&4160017987&583578389575\\
 3&259254948723&571111660440989&47796720049308\\
 4&318278534893&661347985&43971049724\\
 5&260908126900&29917222640283&5659378873682\\
 6&16467340956&328219854&4086161507
\end{array}
\end{equation}
Equations \eqref{appth:coefficients}, \eqref{appth:scalar}, and
\eqref{appth:quotient} provide a finite exact-arithmetic verification of
every entry.

\subsection{Error of the matrix approximation}

\begin{lemma}\label{appth:matrix-error}
For every pair $(n,g)$ in \eqref{appth:integer-table},
\[
 \left\|\widehat E-\exp[-b_0(H_{n,g}+C_n\Id)/2]\right\|_{\rm F}
       <2\cdot10^{-10}.
\]
\end{lemma}

\begin{proof}
Write $d=3^n$. Since $b_0/2=49/8$, the exact target is
\[
 E=e^{-49}\exp(49A)=\sum_{i=0}^{\infty}
                         e^{-49}\frac{49^i}{i!}A^i.
\]
The polynomial with coefficients $q_i=s_i/S_K$ is the truncation of this
Poisson expansion, renormalized to have coefficient sum one. Its coefficient
$\ell^1$ error is twice the omitted Poisson mass, at most $2R$, where
\begin{equation}\label{appth:tail}
R=\frac{49^{193}}{193!}\frac1{1-49/194}<3.152\cdot10^{-33}.
\end{equation}
In bounding the omitted Poisson mass by $R$, we have discarded the
factor $e^{-49}$. This coefficient estimate is used in
\cite[proof of Proposition 1]{FoxGlynn1988}; here it is applied to
powers of a contraction.
Because $A$ is a contraction, the corresponding Frobenius error is at most
$2\sqrt d\,R$.

At each Horner stage the coefficient replacement $q_i\mapsto t_i/Q$
contributes less than $\sqrt d/Q$ in Frobenius norm. The entrywise floor
contributes less than $d/Q$ before undoing the similarity, and hence less
than $\sqrt2\,d/Q$ afterwards, since
$\|J^{-1}\|\,\|J\|\le\sqrt2$. Subsequent propagation is left
multiplication by $A$ and cannot increase the Frobenius norm. Thus the
contraction is used after undoing the similarity; no contraction
property of the generally nonsymmetric integer matrix $D$ is asserted.
The total error is less than
\begin{equation}\label{appth:frobenius-error}
 2\sqrt d\,R+\frac{193(\sqrt d+\sqrt2\,d)}{Q}
 \le54R+\frac{193(27+\sqrt2\,729)}{2^{50}}
 <54R+\frac{432513}{2^{51}}<2\cdot10^{-10}.
\end{equation}
The penultimate inequality uses only $\sqrt2<3/2$; the last one is a
rational comparison using \eqref{appth:tail}.
\end{proof}

Squaring the Frobenius norm converts this estimate into a thermal trace.
In fact, $E$ is self-adjoint and
$\|E\|_{\rm F}^2=\Tr\exp[-b_0(H_{n,g}+C_n\Id)]$. Consequently
\begin{equation}\label{appth:partition-error}
 \left|Z_{n,g}-e^{b_0(C_n-an)}W_{n,g}\right|
 <e^{b_0(C_n-an)}
        \left(4\cdot10^{-10}\sqrt{W_{n,g}}+4\cdot10^{-20}\right).
\end{equation}
This follows from
$|\|E\|_{\rm F}^2-\|\widehat E\|_{\rm F}^2|
\le2\|E-\widehat E\|_{\rm F}\|\widehat E\|_{\rm F}
     +\|E-\widehat E\|_{\rm F}^2$.

\subsection{Rational intervals from the printed table}

The displayed integer floors suffice to verify the needed estimates; no
unprinted digits of $W_{n,g}$ are used. For an entry $k=k_{n,g}$, put
\begin{equation}\label{appth:interval-data}
 L=\frac{k}{10^{16}},\qquad U=\frac{k+1}{10^{16}},\qquad
 r=\frac{\lfloor\sqrt{k+1}\rfloor+1}{10^8},\qquad
 \varepsilon=4\cdot10^{-10}r+4\cdot10^{-20}.
\end{equation}
Then $L\le W_{n,g}<U$ and $\sqrt{W_{n,g}}<r$. In every entry
$L-\varepsilon>0$, so \eqref{appth:partition-error} gives
\begin{equation}\label{appth:partition-interval}
 e^{b_0(C_n-an)}(L-\varepsilon)
       <Z_{n,g}<e^{b_0(C_n-an)}(U+\varepsilon).
\end{equation}
The square root in \eqref{appth:interval-data} is just an integer square
root. To enclose the remaining exponential by rational numbers, set
\[
 T_{80}(y)=\sum_{j=0}^{80}\frac{y^j}{j!}.
\]
For $0\le y\le10$, comparison of consecutive tail terms gives
\begin{equation}\label{appth:scalar-exp}
 T_{80}(y)\le e^y
       \le T_{80}(y)+\frac{2y^{81}}{81!}.
\end{equation}
For a negative argument, take reciprocals and reverse the endpoints.
All arguments $b_0(C_n-an)$ here lie between $-.0544635$ and $9.691073$.
Thus \eqref{appth:interval-data}--\eqref{appth:scalar-exp} require only
rational arithmetic and integer square roots. In particular they give
\begin{equation}\label{appth:Y6-interval}
 .0266319440<Z_{6,1}<.0266319607<.027.
\end{equation}

Using \eqref{appth:partition-interval} and \eqref{appth:scalar-exp} in
\eqref{appth:average} gives the following rational centers, each with
absolute error less than $10^{-7}$:
\begin{equation}\label{appth:moment-centers}
\begin{array}{c|ccccc}
 n&2&3&4&5&6\\\hline
 \text{center of }I(n)
 &.02703074&.01203562&.01193511&.00931560&.00893272.
\end{array}
\end{equation}

Substitution of these centers in the right-hand side of
\eqref{appth:polynomial} gives $.003744031780$. The sum of the absolute
polynomial coefficients is $1.7032$, so the total possible error is less
than $1.7032\cdot10^{-7}$. The right-hand side of
\eqref{appth:polynomial} is consequently less than
\[
 .003744031780+.00000017032=.0037442021<.003745.
\]
Together with \eqref{appth:Y6-interval}, this proves
Proposition~\ref{appth:bounds}.
